\documentclass[11pt]{article}

\usepackage[margin=1in]{geometry}
\usepackage{microtype}
\usepackage{parskip}
\usepackage[hidelinks]{hyperref}
\usepackage{titlesec}
\usepackage{enumitem}
\usepackage{booktabs}
\usepackage{amsmath}
\usepackage{amssymb}

\setlist{nosep, leftmargin=1.5em, topsep=2pt}
\titleformat{\section}{\large\bfseries}{\thesection.}{0.6em}{}
\titleformat{\subsection}{\normalsize\bfseries}{\thesubsection}{0.6em}{}
\titlespacing*{\section}{0pt}{1.6em}{0.5em}
\titlespacing*{\subsection}{0pt}{1.1em}{0.3em}

\hypersetup{
  pdftitle={Before the bank},
  pdfauthor={Robin Langer},
  pdfborder={0 0 0}
}

\title{\textbf{Before the bank}\\[2pt]
\large Why the town could not buy its own output, and the one lever that moved it}
\author{Robin Langer}
\date{7 November 2026}

\begin{document}
\maketitle

\section{The symptom: a town that cannot buy its own output}

The metropolis is a 5$\times$ town --- 75 people, $\sim$40 firms, competing
sellers in every market. Turned on with the ``levelled'' market (incomes
tripled, supply pinched to generation~0, discretionary price discovery), the
nightclubs genuinely sell out and a price moves. But over a five-year run the
pattern is always the same: capacity the firms can clearly afford to build
outruns the town's ability to \emph{buy}, the discretionary price deflates toward
its floor, and the marginal sector folds. The luxury (health-demand) shops are
the canary: they invest early at a high sold-out price, the price then collapses
to~1, and the wage bill of the crew they committed sinks them.

The diagnosis, stated once and never since contradicted: \textbf{the economy is
structurally demand-deficient}. There is plenty of \emph{capacity} and plenty of
\emph{money}; what is missing is purchasing power \emph{reaching the households
that would spend it}.

\section{Six levers, and why each failed}

Each lever below was built additively and measured against the same levelled
town. None cured the deficiency.

\begin{center}
\begin{tabular}{@{}lp{9.2cm}@{}}
\toprule
\textbf{Lever} & \textbf{Result} \\
\midrule
Loosen the price undercut cap & \emph{Worse.} Demand is budget-constrained
(spend $\div$ price), so a higher price buys fewer units and a false ``healthy
price'' lures firms into over-investment. \\
Raise incomes (wage scale $3\!\to\!6$) & \emph{Starvation.} Higher wages
bankrupt the low-margin, inelastic staple producers; \texttt{fed} falls to~0. \\
Population / births & \emph{Collapse.} Poor newborns add labour, not demand; the
swollen idle pool un-caps the invest brake and the clubs over-expand and fold. \\
Generous endowment (a cash gift) & \emph{Transient.} A \emph{stock} depletes in
$\sim$4 months; demand is a \emph{flow} off current cash, so a one-off gift
fades. \\
Bank dividend (wired) & \emph{Inert.} No credit sector, so the bank earns no
profit to distribute. \\
Private lending to firms & \emph{Inert.} The entrepreneurs are not capital-short
--- they hoard 15--25$\times$ the cash a machine needs --- and the idle labour
pool is down to~2, so credit to firms only shuffles idle money between hoarders
and adds supply to a town short of demand. \\
\bottomrule
\end{tabular}
\end{center}

\medskip
The convergent conclusion: \textbf{there is no free demand lever}. Price and wage
are zero-sum redistributions inside a fixed population; a stock gift is not a
flow; and lending to firms is a \emph{supply} channel in an economy whose
binding constraint is \emph{demand}. Sustained purchasing power has to be a
\emph{flow from genuine production into households' hands}.

\section{The diagnosis sharpened: idle wealth is everywhere}

The wealthy are not absent. Rentiers open with 240--600 in cash, owner-operators
150--450, savers up to~360. The question is whether that wealth \emph{acts}.
Tracing the three classic channels through the engine the town actually runs:

\begin{itemize}
\item \textbf{Consumption} fires, but is throttled: a saver spends only a quarter
of current cash, and $\sim$20\% of the town carries the \emph{savings} objective,
which is coded to buy nothing at all --- however rich.
\item \textbf{Real estate} never trades in the monthly loop: households hold one
frozen house each and the market that would move them is never driven.
\item \textbf{Equity and credit} are built but un-driven: the mine$\to$share,
private-loan, and bank-loan machinery all exist, but nothing in the town's loop
ever sets them in motion.
\end{itemize}

So the money does not circulate. Worse than idle rentiers: the \emph{firms} hoard
too, sitting on large payroll reserves. The town is a set of full pools with no
channels between them. That --- not a missing quantity of money --- is what
``demand-deficient'' really names.

\section{Switching on the bank}

The bank is the apparatus that turns an idle \emph{stock} into a household
income \emph{flow}. Switching it on means letting wealth be mobilised rather than
hoarded. Two mechanisms, with very different outcomes so far.

\subsection{Real estate as collateral --- the lever that moved}

An entrepreneur puts up his own house as collateral and takes a bank loan. The
frozen housing stock becomes spendable cash in a consumer's hands, and he spends
a quarter of it. Measured on the levelled town (24 months, equity channel off vs
on):

\begin{center}
\begin{tabular}{@{}lcc@{}}
\toprule
 & \textbf{off} & \textbf{on} \\
\midrule
Luxury shops open & 2 / 6 & \textbf{6 / 6} \\
Luxury capacity (units/mo) & 80 & \textbf{240} \\
\texttt{fed} & 84 & 86 \\
Bank loans taken & 0 & 21 \\
\bottomrule
\end{tabular}
\end{center}

This is the first lever that \emph{created} demand rather than redistributing or
injecting it: mobilising frozen real-estate wealth as credit kept alive the
marginal sector that every other configuration let fold. It is credit-financed
consumption, and demand was the binding constraint, so it bit where supply-side
credit could not. (Caveat: measured at 24 months the one-off injection dominates;
the entrepreneurs now carry debt, and a longer horizon must test whether interest
service claws the gain back.)

\subsection{The dividend --- blocked by the balance sheet}

The intended flow was: the propertied class deposits mines for shares, the bank
earns, and the earnings are paid out as dividends --- firm money reaching
households. The mechanism was built end to end (shares minted, entrepreneurs
mortgaged, interest paid) and it fires correctly --- but the dividend never pays.
The bank's genesis equity is $-(\text{all deposits})$: it is a pure deposit
ledger with no reserve asset, so its equity is pinned at minus the money supply
($\approx -26{,}600$). Interest income raises equity by exactly what is paid, but
a year of mortgage interest ($\approx 1{,}000$) is a rounding error against that
hole, so the dividend gate always refuses. \emph{The dividend is not economically
inert; it is accounting-blocked, and awaits a balance-sheet fix} (deposits backed
by a reserve asset, so genesis equity is $\approx 0$).

\subsection{The completing mechanism: raw materials as income}

The cleanest firm-to-household flow does not need the dividend at all. If a tech
firm must \emph{consume raw materials} to build a machine, and buys them from the
indigenous population who own the mining land, then the money a producing firm
pays for a machine passes through the tech firm and lands as \emph{income} in a
mining household --- which then spends a quarter of it on luxury and
entertainment. Firm hoards, recycled into household demand, via production. This
is the piece being built next; this note will record its measurement when it
lands.

\section{The thesis}

Before the bank, the town has money but no demand, because money does not move:
every classic demand lever is either a zero-sum redistribution, a stock that
depletes, or a supply-side push that a demand-short town cannot use. The bank ---
the credit and capital apparatus --- is what converts idle stocks into household
income flows. The first such flow we switched on (real-estate credit) is also the
first thing that made demand where there was none. The remaining flows (raw-material
income; and, once the balance sheet is fixed, dividends) are the same idea
continued: \emph{put purchasing power in the hands that will spend it, as a flow
from production, not a knob}.

\end{document}
